\subsection{A finite exact search for each prescribed odd prime support}
\label{sec:finite-support-theta}

\paragraph{Finding A19 (an explicit field for all primary exponents).}
Let $C/\mathbf F_Q$ be a smooth projective geometrically connected
genus-two curve in characteristic two. Write $X=C^{(1)}$ for its
one-fold relative Frobenius twist, $J=\operatorname{Pic}^0(X)$,
and $\Theta\subset J$ for its Raynaud theta support. Thus
\[
 \Theta=i_B(X),\qquad
 i_B(R)=\mathcal O_X(R)\otimes B^{-1},
                                                               \tag{FS.1}
\]
where $B$ is the locally exact differential line bundle of degree
one on $X$. This is a genus-two Abel curve translated over
$\mathbf F_Q$; a rational effective divisor for $B$ is not required.
Let $S$ be a finite set of odd primes, and write $J(\overline{\mathbf
F}_Q)\{S\}$ for the points whose orders have every prime divisor in
$S$, with arbitrary exponents. Let $\Phi$ be $Q$-power Frobenius on
$J$, and choose $d\ge1$ such that
\[
            \Phi^d=1\quad\hbox{on }J[\ell]
                         \qquad(\ell\in S).                  \tag{FS.2}
\]
Then
\[
 \Theta(\overline{\mathbf F}_Q)\cap J(\overline{\mathbf F}_Q)\{S\}
   =\Theta(\mathbf F_{Q^d})\cap J(\mathbf F_{Q^d})\{S\}.
                                                               \tag{FS.3}
\]
In particular the left side is finite and is computable by a
terminating finite-field search. One may take
\[
 \begin{split}
 d&=\operatorname{lcm}_{\ell\in S}
                  \operatorname{ord}(\Phi\mid J[\ell]),
       \quad\hbox{or the explicit larger integer}\\
 d&=\operatorname{lcm}_{\ell\in S}
                    |\operatorname{GL}_4(\mathbf F_\ell)|,
 \qquad |\operatorname{GL}_4(\mathbf F_\ell)|
                   =\prod_{j=0}^3(\ell^4-\ell^j).
 \end{split}                                                   \tag{FS.4}
\]
For $S=\varnothing$ take $d=1$ and the torsion group to be $\{0\}$.
No GRH or conjecture about prime distribution is a premise.

\paragraph{Literature scope.}
This is a specialized effective finite-support torsion statement
of Boxall type, proved in full below. The bibliographic locator
for Boxall's general work is
\emph{Sous-vari\'et\'es alg\'ebriques de vari\'et\'es
semi-ab\'eliennes sur un corps fini}, in \emph{Number Theory
(Paris, 1992--1993)}, LMS Lecture Note Ser.~215, Cambridge, 1995,
69--80. The primary paper of Scanlon--Voloch discusses related torsion
finiteness results and lists this reference
\cite[Introduction and reference \textnormal{[Box]}]{ScanlonVoloch1999}.
The original Boxall chapter's theorem text was not accessible in
this check; no theorem number, exact bound, or originality claim
is inferred from its bibliographic record. The proof and the
specific field (FS.4) do not depend on that unavailable text.

\begin{proof}
\emph{1. Two points is the exact geometric intersection bound.}
For every nonzero $T\in J(\overline{\mathbf F}_Q)$,
\[
                  \#\bigl(\Theta\cap(\Theta+T)\bigr)\le2
                                                               \tag{FS.5}
\]
for geometric points. Here is a proof using only genus two.
Suppose three distinct points of this intersection correspond to
pairs $(R_i,S_i)$ on $X$ with
$[R_i-S_i]=T$, for $i=1,2,3$. The $R_i$ are distinct, and so
are the $S_i$, because the Abel map is an embedding.
For $i=2,3$ we have
$R_1+S_i\sim R_i+S_1$. These two effective degree-two divisors
are distinct: equality, with $R_1\ne R_i$, would force
$R_1=S_1$ and hence $T=0$.
Thus the corresponding degree-two linear system moves.
Riemann--Roch in genus two forces its divisor class to be
canonical: for a degree-two divisor $D$, $h^0(D)\ge2$ implies
$h^0(K_X-D)\ge1$, and a degree-zero line bundle with a section
is trivial. Therefore both $S_2$ and $S_3$ are the unique
effective degree-one divisor in the class $K_X-R_1$.
This contradicts $S_2\ne S_3$ and proves (FS.5).

\emph{2. An elementary powering lemma.}
Let $\ell$ be odd, let $V=(\mathbf Q_\ell/\mathbf Z_\ell)^4$,
and let $U\in\operatorname{GL}_4(\mathbf Z_\ell)$ satisfy
$U\equiv1\pmod\ell$. If $v\in V$ is a finite-order point and
$w=(U-1)v$ has exact order $\ell^s$, $s\ge1$, then
\[
        (U^{\ell^{s-1}}-1)v=\ell^{s-1}w\ne0,
                                                               \tag{FS.6}
\]
and this point has exact order $\ell$.
To check the equality put $N=\ell^{s-1}$ and $M=U-1$.
The binomial expansion is
\[
 (U^N-1)v=Nw+\sum_{j=2}^{N}\binom Nj M^{j-1}w.
\]
For $j\ge2$,
\[
 v_\ell\!\left(\binom Nj\right)\ge s-1-v_\ell(j),
 \qquad j-2-v_\ell(j)\ge0.
\]
The first follows from
$\binom Nj=(N/j)\binom{N-1}{j-1}$, and the second uses that
$\ell$ is odd. Since $M\in\ell M_4(\mathbf Z_\ell)$, every
summand after $Nw$ is divisible by $\ell^s$ as an operator on
$w$, and is zero. For $s=1$ the sum is empty. This proves
(FS.6). In addition $U^{\ell^{s-1}}$ fixes every point of
$V[\ell]$, again because $U\equiv1\pmod\ell$.

\emph{3. Isolate one prime component of a nonrational point.}
Let $z\in\Theta(\overline{\mathbf F}_Q)\cap J\{S\}$ and
write $z=\sum_{\ell\in S}z_\ell$ for its primary decomposition.
Put $A=\Phi^d$. Its action on every $T_\ell J$ is congruent
to the identity modulo $\ell$. For a component of order
$\ell^a$, the action of $A$ belongs to the kernel of
\[
 \operatorname{GL}_4(\mathbf Z/\ell^a)
                         \longrightarrow\operatorname{GL}_4(\mathbf F_\ell),
\]
which is an $\ell$-group. Hence the orbit length of $z_\ell$
under $A$ is a power of $\ell$.

Suppose $Az\ne z$ and choose $\ell$ for which
$Az_\ell\ne z_\ell$. For every other prime $t\in S$, choose
a power of $t$ equal to the orbit length of $z_t$ under $A$.
Let $c$ be the product of these orbit lengths. Then
$\gcd(c,\ell)=1$, and $U=A^c$ fixes every other component,
but still moves $z_\ell$: raising a permutation of
$\ell$-power order to the prime-to-$\ell$ power $c$ cannot
turn a nontrivial orbit into a fixed point.

Use the identification
$J[\ell^\infty](\overline{\mathbf F}_Q)
 \simeq T_\ell J\otimes(\mathbf Q_\ell/\mathbf Z_\ell)$.
Let $s\ge1$ be the order exponent of $(U-1)z_\ell$, and put
$\sigma=U^{\ell^{s-1}}$. By (FS.6),
\[
                      \sigma z=z+T,
 \qquad T\in J[\ell]\setminus\{0\},\qquad \sigma T=T.
                                                               \tag{FS.7}
\]
The other primary components remain fixed. The map $\sigma$
on geometric points is a positive integral power of $Q$-power
Frobenius. As $\Theta$ is defined over $\mathbf F_Q$, all
\[
                \sigma^jz=z+jT\qquad(0\le j<\ell)
\]
belong to $\Theta$. They are $\ell$ distinct points. Each also
belongs to $\Theta+T$, by using the preceding point in the
cyclic translation orbit. This contradicts (FS.5), since
$\ell\ge3$. Consequently $Az=z$, which is exactly
$z\in J(\mathbf F_{Q^d})$. The reverse inclusion in (FS.3)
is immediate.

\emph{4. An explicit order bound and terminating computation.}
Let
\[
 h_d=\#J(\mathbf F_{Q^d}),\qquad
 N_S=\prod_{\ell\in S}\ell^{v_\ell(h_d)}.
                                                               \tag{FS.8}
\]
Every point in (FS.3) has order dividing $N_S$, and
\[
                 N_S\le h_d\le(Q^{d/2}+1)^4.                 \tag{FS.9}
\]
The last inequality is the genus-two Jacobian Weil bound
\cite[Theorem~19.1(c)]{MilneAV}.
If $P_J(T)=\det(T-\Phi)$ is its integer Weil polynomial, then
\[
 h_d=\prod_{j=1}^4(1-\alpha_j^d)
       =\operatorname{Res}_T(P_J(T),T^d-1),                 \tag{FS.10}
\]
where the $\alpha_j$ are its roots; the sign is positive
because there are four roots
\cite[Theorem~19.1(b)]{MilneAV}.
Thus $h_d$, its valuations at the finitely many primes in
$S$, and $N_S$ are exact finite integer computations.
Using the explicit second choice in (FS.4) does not require
computing the torsion representation first.

Enumerate $X(\mathbf F_{Q^d})$, form the line-bundle classes
$i_B(R)$ in the finite Jacobian, and keep precisely those killed
by $N_S$. This enumerates the entire geometric intersection
in (FS.3). Their exact orders, and any desired doubling-orbit
coverage within their cyclic subgroups, can then be tested by
finite group operations. A smaller bound may be obtained by
replacing $h_d$ with the exponent of its $S$-primary subgroup,
but this is unnecessary for termination. No assertion of
practical running time at the potentially large degree $d$ is made.
\end{proof}

\paragraph{Consequence for the remaining composite-order search.}
For each fixed prime $q$ in A17, taking $S=\{3,q\}$ now makes
the search over \emph{all} orders $3^a q^b$ finite, not only
the search at order $3q$. The actual order must still satisfy
the unit-group test or the same-subgroup orbit-coverage condition.
The primes $q$ range over an unbounded set, so this result does
not prove that one of the finite searches succeeds for every
curve.

There is also an exact reformulation of the squarefree gate.
For $q\ne3$, the curve $[3](\Theta)$ contains a nonzero point
of order $q$ if and only if $\Theta$ contains a point of order
$q$ or $3q$. Indeed a lift $z\in\Theta$ of such a point
satisfies $3qz=0$ but $3z\ne0$, so has one of those two orders;
the converse follows by multiplying by three.
For $q$ in A17 either order meets the genus-two sufficient
condition of A18. Applying Poonen's multiples theorem to the
generating curve $[3](\Theta)$ still permits a varying additional
integer multiplier. It does not put a prime-order point on that
curve itself. Thus neither this reformulation nor (FS.3)
replaces the remaining incidence assertion by prime density.
